\chapter{Conclusion and Outlook}\label{ch:conclusion}

This chapter takes stock: for each research question RQ0--RQ4 (Section~\ref{sec:rqs}) it puts the
sub-questions first, names the state of their answer in three grades --- answered, structurally
answered, open ---, traces per question the chain of evidence through the chapters and its
consequences, assigns the hypotheses to the target path, names the limits of the system from the
user's point of view, and sketches the next steps.

\section{Answering the Research Questions}\label{sec:answers}
The answers follow the scheme claim, evidence, validation~\cite{shaw2003writing}: a question counts as
\emph{answered} when a test or a measurement substantiates it; as \emph{structurally answered} when
the apparatus that delivers the answer is implemented and tested but the measurement is not
available; as \emph{open} when neither structure nor measurement is available. Three kinds of
evidence occur: the structure (a code object that supports the answer), the test (a static check or a
contract test that secures it), and the measurement (a number together with its run); for the
question about the state of the art of the workloads (RQ2.d) the literature evidence is added. A
\enquote{structurally answered} does not replace the number: without measured values the main question
remains unanswered even if its structure is sound. Table~\ref{tab:ff-beweiskette} lists status,
evidence, and remainder per main and sub-question; the form of reporting follows the rules of empirical
software and systems research~\cite{wohlin2012experimentation,runeson2009guidelines} and the demand to
state the reference quantity, the platform, and the values not collected at every
result~\cite{blackburn2016truth,hoefler2015benchmarking}.

Orthogonal to all five questions, the introduction posed a categorization question as a question:
whether the axes of the design space separate cleanly into organ, system, and measurement axes ---
and how such a separation is to be accomplished where a design does not provide for it. Both are
answered (Section~\ref{sec:axis-triad}): the axes separate along the three concerns into three
\emph{realms} --- organ, system, measurement --- with one class root each, one offering catalogue each
(three generated registries: the organ-axis registry of the cache engine, the system-axis registry,
and the measurement-axis registry), and one stamp line each at the binary
(Section~\ref{sec:stamp-model}); the three line carriers of the fingerprint --- organ, system, and
measurement line --- are, as types, pairwise unconstructible from one another, so that a swap does
not compile. Only the organ main axes enter the \texttt{binary\_id}; system and measurement axes are
declared binary-id-neutral in their registries. The separation is held by compile-time
\emph{re-entry guards} (Section~\ref{sec:impl-system-axes}): \texttt{static\_assert} checks reject
compiler, scheduling, and load framework as system main axes --- the scheduling is a sub-axis of the
target ISA \texttt{target\_isa} with five dimensions, the compiler a sub-axis group of the complex
bracket \texttt{build\_target\_complex}, the load framework the meta-meta axis of the measurement
realm ---, and the flag grammar of the version stamps admits only catalogued flag chains. This triad
is the second shape of the separability problem and supports the answers to the following
sub-questions.

\paragraph{RQ0 (Main question).} The question is by how much active, measurement-driven cache-aware
decisions --- page type, memory layout, value placement, prefetching --- outperform their passive,
statically dimensioned counterparts on different CPU platforms; the extension RQ0.E asks whether a
measured gain can be attributed unambiguously to the organ, the system, or the measurement concern.
The \emph{status}: RQ0 is open. Without the end-to-end runs per platform there is no number; the
structure that delivers the number is implemented. The same axis-constant structure is built once
actively --- with a runtime-dynamic cache-line policy (the policy selector derives the cache-line
adaptation from an evaluated profile aggregate of operation mix and working set) and
measurement-driven selection of the best binary --- and once passively --- with a statically fixed
organ choice of memory layout, prefetch, value placement, and node type as segments of the
\texttt{binary\_id} --- and the two are compared under an identical compiler and flag basis
(default \texttt{-O2}; \texttt{-O3} only as an opt-in under warning and with its own toolchain
identity). The measurement machines surveyed are the three registered machine signatures: prod1
(AMD Zen~5, Ryzen~9~9950X3D), prod2 (Intel Core i9-12900K, Alder Lake --- a hybrid CPU with % chktex 29
P- and E-cores, without AVX-512), and an ODROID board with Intel Gracemont cores; the server
platform Sapphire Rapids named in the question is provided for in the platform detection (ZIH node
Barnard, Xeon~8470) but not measured (n/a; Section~\ref{sec:eval-platforms}). RQ0.E is structurally
answered, and the separation it requires is the triad: an organ effect permutes the
\texttt{binary\_id}, a platform concern stands at the binary as system stamp line and build version
identifier without permuting it, and a measurement concern remains binary-id-neutral --- no
measurement instrument changes \emph{what} is built. The \emph{evidence} is
Section~\ref{sec:axis-triad} (triad; structure), Section~\ref{sec:impl-system-axes} (re-entry
guards; test), and Section~\ref{sec:eval-method} (measurement plan; structure). The
\emph{consequence}: a quantification presupposes the full build and the run of series~A
(Section~\ref{sec:outlook}); the effect of cache-line awareness from the synthesized measurement
curves is the first evaluation step after the first measurement run
(Section~\ref{sec:heuristics-curves}).

\paragraph{RQ1 (Composition).} The question is how a search-algorithm design can be canonically
decomposed into orthogonal axes and systematized in a permutation matrix, such that
state-of-the-art methods become reconstructible bias-free as profiles, ideally from their original
code, and arbitrarily permutable; RQ1.a asks about the category cut of the decomposition, RQ1.b
about the organ couplings and the behaviour under a solely varied ISA/SIMD axis. The \emph{status}:
RQ1 is structurally answered, RQ1.a answered, RQ1.b open in its empirical part. The canonical
decomposition carries: \emph{eighteen} organ main axes in fixed slot order T0--T17 --- as the
constant name array of the axis-path serialization they are code, not convention --- span the
permutation matrix (Section~\ref{sec:axis-triad}); to these is added an optional
\emph{organ meta-meta axis} for disk input/output \texttt{disk\_io}, carried in the organ registry
as the nineteenth entry (a compile-time stage with the runtime option \texttt{staging},
binary-id-neutral), but not selectable via experiment profile, because the profile parser does not
read its registry block; the profile connection is planned but not implemented. It is the thesis's
example of the meta-meta level M3, the indirection of a description: it stamps in the bracket
appendage of the organ line, and the bracket depth of an entry is its meta-meta level
(Section~\ref{sec:stamp-model}). Two reference quantities are to be kept apart: the golden
reference catalogue varies seventeen axes over two values each and pins the persistence target to
one value --- $2^{17} = 131072$ binary identities per system permutation as a regression detector
with a fixed CRC-64 checksum\footnote{CRC-64: a cyclic redundancy check with a 64-bit checksum
(here the variant ECMA-182), computed over all identities in canonical order; a changed identity
changes the checksum.} over all identities (Section~\ref{sec:impl-contracts}) ---, whereas the
cardinality of a planned space is not a constant but the planner's computation per profile and
platform (subcommand \texttt{simulate}); all eighteen axes are planned for variation (not
measured). The 33 corpus works P01--P33 each possess a profile XML, among them 30 full profiles of
standalone search designs (eight Rank-1 and 22 Rank-2/3) and three abstract profiles (OLC, LOUDS,
VAMPIR) that occupy only individual axes (Section~\ref{sec:sota-instances}); the registry of the
paper probes holds this translation at compile time against duplicates and gaps. The variadic
mapping --- one template parameter for the \texttt{std::vector} API, two for \texttt{std::map},
$N>2$ for \texttt{std::map<K,\,std::tuple<\ldots>>} (Appendix~\ref{app:interfaces}) --- makes
algorithm- as well as container-oriented patterns expressible on one engine, and the family triad
Map, Container, Graph orders the genera beneath it (Section~\ref{sec:anatomy-levels}). RQ1.a is
answered in the affirmative: the decomposition remains orthogonal only because platform and
measurement concerns are kept out of the permutation matrix, and the re-entry guards hold this cut
(Section~\ref{sec:impl-system-axes}). On RQ1.b: the coupling node type $\times$ SIMD $\times$ path
compression named by the question falls apart --- node type $\times$ path compression remains a
not-fully-orthogonal organ coupling that is marked as such and restricted during permutation,
whereas SIMD is not an organ axis but a build control quantity of the system axis
\texttt{external\_utils} (a runtime sub-axis with the options \texttt{no\_extension},
\texttt{avx2}, and \texttt{avx512} and a feature catalogue of 23 flags, binary-id-neutral): one and
the same \texttt{binary\_id} is built under several system permutations and compared with an
otherwise unchanged structure; the empirical answer on cache-update and cache-line fill behaviour
is a measurement and is not available. The \emph{evidence} is Section~\ref{sec:axis-triad}
(decomposition and realms; structure), Section~\ref{sec:mess-chain} (catalogue and reference
quantity; structure), and Section~\ref{sec:impl-contracts} (catalogue checks and golden checksum;
test). The \emph{consequence}: RQ1 structurally answered; RQ1.b (fill behaviour) presupposes the
full build.

\paragraph{RQ2 (Measurement methodology).} The question is how a reproducible measurement pipeline
with hardware-counter collection must be designed so that every axis permutation is measured
fairly --- under an identical compiler and flag basis --- in three measurement series over three
granularities, and under which platform-calibrated criteria a platform counts as verified; RQ2.a
asks whether the measurement instruments themselves must be axes of their own, RQ2.b about the
re-measurement of contradictory literature findings, RQ2.c about extensibility, RQ2.d about the
workloads that belong to the state of the art specifically for search algorithms. The
\emph{status}: RQ2 is structurally answered, RQ2.a and RQ2.d answered. Measurement is split into a
production mode and an experiment mode, switched via the master switch of the measurement
compilation (\texttt{COMDARE\_\allowbreak{}MEASUREMENT\_\allowbreak{}ON}): the experiment mode
always measures --- a release build without measurement sensors is a separate, deliverable
artefact that forces the switch off, not a measurement run that has been switched off ---, and in
the store only what is missing there is re-measured; the compile-time-static key schema of the four
stocks is the design of this recognition, the full build-out of the store is not implemented
(Section~\ref{sec:lager}). The \texttt{ExperimentDriver} controls the sequence host-side: the
runner loads the modules and runs the phases load, execute, measure, and teardown, the driver runs
execution and measurement per module and collects the results in the \texttt{ResultAggregator}.
Every series is collected on three \emph{measurement levels}, in exactly this order: the wall-clock
level (\texttt{wallclock}; the overall runs in the CacheEngineBuilder over the load frameworks),
the macro level (\texttt{macro}; the family and genus interface operations of the tier binary), and
the micro level (\texttt{micro}; axis interface and success parameters); w/ma/mi is the registered
alias, and the three levels are the measurement-tooling main axis, which per choice is compiled
into one CacheEngineBuilder route (Section~\ref{sec:mess-chain}). RQ2.a is answered in the
affirmative: the measurement instruments are a realm of their own in the planner --- the load
framework as meta-meta axis, the measurement tooling as main axis, the measurement categories as
sub-axis (sixteen building blocks in the registry; planned is the collection per axis-interface
parameter set, whose catalogue per axis is not implemented) ---, always present as soon as
measurement takes place, and binary-id-neutral. The separation is accomplished by the double root
cut: the system realm is rooted at the CacheEngineBuilder, the measurement realm at the planner, and
the measurement axis drives, via the system-axis root, the release of the measurement source in the
form of the CacheEngineBuilder type --- the legend chain measurement axis $\to$ CacheEngineBuilder
type $\to$ system axes $\to$ tier binary is code as a naming scheme
(Section~\ref{sec:axis-triad}). For the workload, the driver passes a default workload through per
module, and a profile brings its own; planned is the full permutation over all available workloads,
so that no design wins only under its self-chosen load profile --- the apparatus counters the
\emph{load-profile selection bias}\footnote{Load-profile selection bias (workload selection bias):
the distortion of a comparison that arises because the load profiles are chosen such that one
design performs preferentially well under them; the benchmarking literature therefore demands a
broad, justified set of load profiles~\cite{blackburn2006dacapo,blackburn2016truth}.} with the full
matrix of all load profiles against all living-being binaries, because even inconspicuous setup
decisions systematically distort measurement results~\cite{mytkowicz2009wrongdata}; implementation
status: the full permutation is not implemented.

On RQ2.b (organ axes): the contradictory literature findings on the optimal node/cache-line size
--- one cache line per node for CSS and CSB\textsuperscript{+} trees~\cite{rao1999css,rao2000csb}
versus sixteen and more cache lines for Hankins and Patel~\cite{hankins2003nodesize} (nodes of
512 bytes and larger at 32-byte cache lines) as well as several per node for the prefetching
B\textsuperscript{+}-trees of Chen, Gibbons, and Mowry~\cite{chen2001prefetch} --- are not adopted
from their original setups but re-measured axis-isolated under otherwise equal conditions: node
type, memory layout, and prefetch are organ axes that are varied individually while system and
measurement axes stay constant; the number is a measurement and is not available. On RQ2.c
(extensibility): a new probe joins as a living being of its family into the genus --- family is the
first level (Map, Container, Graph), genus the second, and the genus build admission carries the
slot sets per family as a compile-time verdict ---; it docks at the probe dock, passes the
three-stage evaluation (Section~\ref{sec:mess-chain}), and its merge with the cache engine bears
one of the registered names: \texttt{Verbund1\_CeOnly} (no probe directive),
\texttt{Verbund2\_Replace} (the probe axis replaces the cache-engine axis, with fallback),
\texttt{Verbund2\_Hybrid} (cache engine and probe hybrid per probe; declared, not materialized ---
the validator blocks the token \texttt{merge} as long as the strategy value is not implemented), or
\texttt{Verbund3\_Union} (union of all variants; prepared interface,
Section~\ref{sec:eval-series}). At the binary the merge remains distinguishable via the
hierarchically extended algorithm names and flags of the organ line --- a probe organ carries its
namespace prefix in the entry ---; a stamp line of its own is not needed, and a version suffix does
not denote the mixture: the flag \texttt{e} denotes an efficiency-core algorithm
(Section~\ref{sec:stamp-model}). Whether the probe matches an existing family, founds a new one, or
constitutes a hybrid form is decided by the family cut (Section~\ref{sec:anatomy-levels}); a family
of its own would arise only with a fundamentally different base axis set. The operating mode asked
for in RQ2.c, which offers the design constituents of exactly one publication as a closed axis set
in isolation, is not implemented and remains a subject of the outlook
(Section~\ref{sec:outlook}). On RQ2.d (workload state of the art): the state of the art of the
workloads for search algorithms is formed by the key-value convention YCSB with its six load
profiles A--F (each an operation mix over a key access distribution, from 50/50 read/update
Zipf-distributed to short range scans and read-modify-write), the index benchmark SOSD, the TPC
family (TPC-C, TPC-DS), the integration into real systems (SuRF in RocksDB), real string corpora
(URLs, DNA, dictionary terms), as well as SPEC~CPU and CloudSuite in the hardware-near works
(Section~\ref{sec:workload-basics-known}); these frameworks generate loads that are not directly
comparable with one another, and each publication typically chooses the load profile under which
its own design performs best. From these the thesis adopts YCSB as the registered load framework
(the meta-meta axis \texttt{load\_framework} of the measurement realm) with the runtime dimension
\texttt{workload} over the six tokens \texttt{YCSB\_A} to \texttt{YCSB\_F}, routed per permutation
from the traversal tag of the profile or its \texttt{<expected\_workload>} tag, and runs them over
real string datasets with a reproducibility record (\texttt{url}, \texttt{protein},
\texttt{tpcds-id}, \texttt{trec-terms}, \texttt{dna}, \texttt{xml};
Section~\ref{sec:eval-workloads}), where \texttt{xml} possesses only a specification without
checksum because the file is not available locally (n/a); the remaining conventions are not
registered as load frameworks (n/a). Implementation status: for profile-less modules the driver
falls back to exactly one workload, the workload option of the run; the full permutation over all
load profiles is planned but not implemented (Section~\ref{sec:eval-workloads}). On platform
verification: a platform counts as verified when the conformance gate confirms, per loaded tier
binary, the seventeen \texttt{std::map} operations marked in Appendix~\ref{app:interfaces} --- the
operation set of its twelve core check blocks, among them \texttt{at}, \texttt{contains},
\texttt{count}, \texttt{insert\_or\_assign}, \texttt{erase}, \texttt{clear}, \texttt{lower\_bound},
and \texttt{upper\_bound} --- against \texttt{std::map} as oracle and the platform probe covers the
reported axis subset; the five drive operations of the module boundary (\texttt{insert},
\texttt{lookup}, \texttt{erase}, \texttt{clear}, \texttt{size}) are the ABI primitives from which
the host composes the remaining operations, not the scope of the check, and the full build-out to
the complete \texttt{std::map} interface is not implemented; the platform-calibrated thresholds are
not available, their determination presupposes the measurement runs. The \emph{evidence} is
Section~\ref{sec:mess-chain} (three-stage evaluation and merge; structure),
Section~\ref{sec:measurement-system} (measurement levels and realm; structure),
Section~\ref{sec:anatomy-levels} (family cut; structure), Appendix~\ref{app:interfaces} (contract
subset; test), and, for RQ2.d, Section~\ref{sec:workload-basics-known} (workload conventions;
literature) and Section~\ref{sec:eval-workloads} (workload routing and records; structure). The
\emph{consequence}: RQ2 structurally answered, RQ2.c conceptually answered, RQ2.d answered; the
isolation mode, the number for RQ2.b, and the full permutation of the workloads are open.

\paragraph{RQ3 (Probe comparison).} The question is whether the probe PRT-ART wins against the
eight Rank-1 designs --- ART~\cite{leis2013art}, HOT~\cite{binna2018hot}, Masstree, CoCo-Trie,
START, B\textsuperscript{2}-tree, Wormhole, SuRF --- and the Rank-2/3 SOTA \emph{not} through
asymptotic novelty but through better microarchitectural properties --- higher cache-line
utilization, fewer LLC and dTLB misses, lower branch cost, smaller hot search paths --- for exact
lookup, prefix lookup, and prefix enumeration, measured in the mandatory measurement series A, B,
and C\@; RQ3.a asks about the cheapest local page representation per branch point, RQ3.b about
the compile-time force. The \emph{status}: RQ3 is open, RQ3.b answered, RQ3.a structurally
answered (prescription without measurement evidence). The comparison is not available and
presupposes the first measurement runs; the apparatus for it is implemented: the PRT-ART profile,
the three mandatory measurement series --- the series is derived mechanically from the merge stage
(\texttt{Verbund1} and \texttt{Verbund2} yield series~A, \texttt{Verbund3} series~B, series~C is
cross-build; Table~\ref{tab:stage-series}) ---, the \texttt{ExperimentDriver} with auto-pickup of
the 33 profile XMLs from the profile directory, and the ABI-conformant adapter with the
\texttt{std::map} core contract as a host-side composition layer over the exact-lookup primitive
of the module boundary, flanked by insertion and deletion, with ordered iteration as an additive
scan capability (Appendix~\ref{app:interfaces}). The conformance gate checks the seventeen
operations marked in Appendix~\ref{app:interfaces} in twelve core check blocks against
\texttt{std::map} as oracle; the full build-out to the complete interface is planned but not
implemented. The notification hook of the measurement layer is a single-slot hook per axis, not
compiled in by default (opt-in), and not part of the contract. On RQ3.a: array-indexed page --- the
direct-index pages Array$_{256}$ and Array$_{65535}$ --- versus vector page --- the vector pages
Vector$\langle$u8,u8$\rangle$ and Vector$\langle$u16,u16$\rangle$ --- per branch point and the
switching rules by occupancy density and access type are part of the PRT-ART profile; the probe
switches at fixed density thresholds of 25, 50, and 75~percent (Section~\ref{sec:prtart-demo}),
and the platform-calibrated rules are prescriptions without measurement evidence as long as the
runs do not substantiate them. RQ3.b is answered in the affirmative and has become the
construction principle: every axis divides into a compile-time-fixed \emph{main axis} and
runtime-variable \emph{sub-axes} --- the main axes are compiled in statically, for organ axes as a
\texttt{binary\_id} segment, for system axes as stamp together with the build version identifier,
while sub-axes with the stage \texttt{runtime} vary their parameters without recompilation; the
organ registry carries the stage per sub-axis, the system registry the sub-axis group
\texttt{build\_toolchain} of the complex bracket (Section~\ref{sec:measurement-system}). The
\emph{evidence} is Section~\ref{sec:prtart-demo} (probe and page representation; structure),
Section~\ref{sec:sota-instances} (33 profiles; structure), Section~\ref{sec:measurement-system}
(main and sub-axis; structure), and Appendix~\ref{app:interfaces} (core contract; test). The
\emph{consequence}: the number presupposes series~A\@; the hypotheses H1--H3 supply the evaluation
criteria.

\paragraph{Hypotheses and target path.} The evaluation hypotheses H1--H3
(Section~\ref{sec:eval-hypotheses}) were formulated before the apparatus; the target path developed
across the chapters --- triad of the realms, stamp identity, measurement levels --- has not laid
them as a foundation but ranks them as touchstones of the evaluation: H1 (page-type cost) and H3
(distribution of the value placement) are checked in series~A per load profile, H2 (code quality
per source) via the tool-computed quality record. Its data basis is the weighted
\texttt{cppcheck} finding density per thousand physical lines over the vendored original sources;
the record carries 33 entries, 22 of them as n/a because a reconstruction possesses no vendored
original source, and PRT-ART possesses no original repository (n/a). The hypothesis check follows
the requirement to name criterion and threshold before the measurement~\cite{wohlin2012experimentation},
and the number of repetitions per cell follows the required confidence~\cite{kalibera2013rigorous}.
H2 counts as accepted if, over the full profiles with a vendored original source, the Spearman rank
correlation between finding density and measured throughput per load profile and cell (platform,
workload, regime) is negative with $|\rho| \geq 0.5$ at $\alpha = 0.05$; as rejected at
$|\rho| < 0.3$; in between as undecided; below eight profiles with a score H2 is not testable. The
manual audit over seven assessment axes remains the qualitative classification. The PRT-ART
telemetry metrics H1--H3 are measurement organs of the probe, not these hypotheses
(Section~\ref{sec:eval-hypotheses}); and the variant ChainRef of the value placement is the subject
of H3, it is no abort criterion for H2.

\paragraph{RQ4 (Compile time vs.\ run time).} The question is which axis properties must be fixed
at compile time to produce platform-optimized, provenance-verifiable binaries per permutation, and
which can or must be selected at run time (bit-encoded binary identity per permutation,
runtime-adaptive layout selection) without violating comparability (identical compiler and flag
basis) --- and whether the measurement makes it possible to generate ISA+OS-optimized binaries for
a search-algorithm recombination whose heuristic optimization is perfectly tailored to its runtime
environment; RQ4.a asks whether the entire sequence --- from the platform-specific build of all
binary permutations via the measurement run and the selection of the best measured binary or
binary recombination up to the evaluation in CSV/xlsx files, LaTeX, and diagrams --- can be
generated end to end from one XML description of the experiment; the extension RQ4.E asks whether
the boundary runs along the axis categories.
The \emph{status}: RQ4 is structurally answered, RQ4.a structurally answered, RQ4.E answered. The
boundary runs along the categories: the three system main axes \texttt{target\_isa},
\texttt{operating\_system}, and \texttt{external\_utils} together with their complex bracket
\texttt{build\_target\_complex} (Section~\ref{sec:complex-axis}; no fourth main axis) belong to the
compile-time identity of the binary; they appear as system stamp line and build version
identifier, never as a \texttt{binary\_id} segment. Provenance is secured by the tier fingerprint:
a SHA-512 over eleven members (fingerprint format~6), among them the overlay member as a SHA-512
over the concatenated source files of the overlay cut --- per axis the file set of its own
implementation plus the shared anatomy substance ---, so that a pure source-text change yields a
new identity, and the toolchain member with the build-compiler identifier and the optimization
identity; the planner fingerprint is a SHA-256 over its own preimage. The query
\texttt{get\_compiler()} of an axis, by contrast, returns the original compiler of the publication
from the manifest of the vendored source --- a statement of origin, not a build identifier.
Measurement choices remain runtime-variable; layout and page selection remain selectable via the
mixed-radix order of the \texttt{binary\_id} (the static binary view) and the runtime-dynamic
cache-line policy, without violating the common compiler and flag basis. Where an axis needs both
sides of the boundary, it is two-way split: load framework, telemetry, and hardware detection run
across the runtime of the planner \emph{and} the compile time of the builder --- the sub-axis
release of the earlier stage becomes the main-axis assumption of the later stage
(Section~\ref{sec:axis-triad}) ---, and the hardware detection sets the compile-time stamps of the
tier binaries from the actual properties acquired at builder runtime (Section~\ref{sec:hw-probe}).
What is stamped is always the platform \emph{class}, never a measured value; the \emph{fit stamp},
which checks the stamped class on every measurement run against the runtime verdict of the
hardware detection, is designed but not implemented (Section~\ref{sec:outlook}). The
forward-looking sub-question of whether the measurement itself can tailor ISA- and
operating-system-optimized binaries for a recombination is taken up by the outlook as heuristics
automation; the operating-system dimension is provided for with three declared values of the
system axis (\texttt{linux}, \texttt{windows}, \texttt{macos}, each with its own probe) and the two
measurement regimes in the measurement plan --- the second regime, an immutable operating system,
is designed, not implemented (Section~\ref{sec:eval-platforms}). On RQ4.a (declarative chain): the
chain is laid out declaratively end to end and built as two binaries. The planner
\texttt{comdare-experiment-planner} checks the user XML read-only against the compiled-in registry
catalogues and emits the deterministic plan from the one director walk --- two runs are
byte-identical, the plan header annotates the registry trio, and every build step reports the
actually generated file stem from the same orchestrator function that the build also calls ---;
the measurement driver \texttt{comdare-messung-driver} emits the system permutations and tier
build orders of its released build space and runs the seven-phase measurement pipeline from
\textsc{Enumerate} to \textsc{Persist} (CSV and JSON); the downstream nine-stage toolchain
renders from that the xlsx workbook, the booktabs tables, the TikZ diagrams, and the appendix
files that the manuscript includes --- one XML source, one build channel, no step by hand
(Section~\ref{sec:impl-pipeline}, Figure~\ref{fig:eval-xml-kette}). The selection of the best
binary or recombination is registered and validatable as the work mode \texttt{compare} of the run
chain; its mode-specific sequence and the \texttt{release} are planned but not implemented ---
both are selectable, validatable labels with \texttt{release}-equivalent build semantics
(Section~\ref{sec:anatomy-levels}).
Implementation status: the chain is demonstrated in a first end-to-end run --- the smoke
measurement series over 43 permutations in the appendix is its product ---; the headless service,
the XML creation mode, the \texttt{<publish>} section, and the CSV tree form are planned but not
implemented. The \emph{evidence} is Section~\ref{sec:complex-axis} (complex bracket; structure),
Section~\ref{sec:hw-probe} (two-way split detection; structure), Section~\ref{sec:stamp-model}
(stamp lines; structure), Section~\ref{sec:impl-system-axes} (re-entry guards and order checks;
test), and, for RQ4.a, Section~\ref{sec:impl-pipeline} (planner, driver, and measurement pipeline;
structure) and Section~\ref{sec:mess-chain} (chain from the XML to the tier binary; structure).
The \emph{consequence}: RQ4 and RQ4.a structurally answered; heuristics automation, fit stamp, and
the sequence of \texttt{compare} and \texttt{release} open.

\begin{table}[htbp]\centering\footnotesize
\begin{tabular}{@{}>{\raggedright\arraybackslash}p{1.3cm}>{\raggedright\arraybackslash}p{3.0cm}>{\raggedright\arraybackslash}p{2.0cm}>{\raggedright\arraybackslash}p{4.4cm}>{\raggedright\arraybackslash}p{3.3cm}@{}}
\toprule
Question & Sub-question & Status & Evidence (sections; kind) & Remainder \\
\midrule
RQ0 & main question: active, measurement-driven versus passive, static & open & {\ref{sec:axis-triad}} (structure), {\ref{sec:impl-system-axes}} (test), {\ref{sec:eval-method}} (measurement plan) & number from series~A after the full build \\
RQ0.E & attributability of a gain to organ, system, measurement & structurally answered & {\ref{sec:axis-triad}}, {\ref{sec:stamp-model}} (structure) & measurement that exercises the attribution \\
RQ1 & canonical decomposition, permutation matrix, profiles & structurally answered & {\ref{sec:axis-triad}}, {\ref{sec:mess-chain}} (structure), {\ref{sec:impl-contracts}} (test) & run over all eighteen axes; profile connection of the meta-meta axis \\
RQ1.a & category cut of the decomposition & answered & {\ref{sec:axis-triad}} (structure), {\ref{sec:impl-system-axes}} (test) & --- \\
RQ1.b & organ couplings, ISA/SIMD variation & open (measurement) & {\ref{sec:axis-triad}} (structure) & cache-line fill behaviour per system permutation \\
RQ2 & measurement pipeline, three series, three levels, verification & structurally answered & {\ref{sec:mess-chain}}, {\ref{sec:measurement-system}} (structure) & full permutation of the workloads; calibrated thresholds \\
RQ2.a & measurement axes as a realm of their own & answered & {\ref{sec:axis-triad}} (structure), {\ref{sec:impl-system-axes}} (test) & catalogue of the success parameters per axis \\
RQ2.b & node/cache-line size axis-isolated & structurally answered & {\ref{sec:measurement-system}} (structure) & measurement over node type, layout, prefetch \\
RQ2.c & extensibility, integration of new probes & answered (conceptually) & {\ref{sec:anatomy-levels}}, {\ref{sec:mess-chain}} (structure), {\ref{app:interfaces}} (test) & paper-isolation mode; materialization of \texttt{Verbund2\_Hybrid} \\
RQ2.d & workload state of the art & answered & {\ref{sec:workload-basics-known}} (literature), {\ref{sec:eval-workloads}} (structure) & full permutation of the workloads; only YCSB registered as load framework \\
RQ3 & PRT-ART against SOTA in series A--C & open & {\ref{sec:prtart-demo}}, {\ref{sec:sota-instances}} (structure) & series~A\@; full build-out of the contract \\
RQ3.a & page representation, switching rules & structurally answered (prescription without measurement evidence) & {\ref{sec:prtart-demo}} (structure) & platform-calibrated rules from the runs \\
RQ3.b & compile-time force & answered & {\ref{sec:measurement-system}} (structure), {\ref{sec:impl-contracts}} (test) & --- \\
RQ4 & compile time versus run time, provenance & structurally answered & {\ref{sec:stamp-model}}, {\ref{sec:hw-probe}} (structure), {\ref{sec:impl-system-axes}} (test) & fit stamp; heuristics automation \\
RQ4.a & declarative chain from the XML to the evaluation & structurally answered & {\ref{sec:impl-pipeline}}, {\ref{sec:mess-chain}} (structure) & sequence of \texttt{compare} and \texttt{release}; service, creation mode, CSV tree \\
RQ4.E & boundary along the axis categories & answered & {\ref{sec:complex-axis}}, {\ref{sec:stamp-model}} (structure) & --- \\
\bottomrule
\end{tabular}
\caption{The chain of evidence per main and sub-question: status in three grades (answered,
structurally answered, open), the supporting sections with their kind of evidence (structure, test,
measurement; for RQ2.d the literature evidence), and the open remainder without a
number.}\label{tab:ff-beweiskette}
\end{table}

\paragraph{Implementation status.} The declared build scope --- per system permutation the 131072
organ permutations of the golden reference catalogue; the catalogue comment estimates roughly
224~GB and roughly 24~hours per complete build, an estimate, not a measured value --- has not been
carried out; the end-to-end runs of the measurement driver across the target platforms (x86-64
with AVX2 and AVX-512, aarch64 with NEON and SVE\footnote{NEON and SVE\@: the SIMD extensions of
the ARM architecture --- NEON with a fixed 128-bit width, SVE (Scalable Vector Extension) with an
implementation-dependent vector width.}) are not available, nor is the delivery measurement over
the 320-identity catalogue. Substantiated by a run are two partial executions: the first
end-to-end run of the chain from the experiment XML via build and measurement to the included
appendix (Section~\ref{sec:eval-xml}) --- a measurement job of the smoke stage on the AMD Zen~5
machine\footnote{Provenance of the smoke run: pipeline 15800, measurement job 377503, host prod1
(Appendix~\ref{app:measurements}); a preceding first end-to-end run of the same chain (pipeline
15764, job 376333) yielded 1199.047~ns per operation at 10000 operations, its measuring machine is
not documented and the difference not assessed.} with 10000 operations and 671.5~ns per operation
(p50 of the \texttt{lookup} operation 750~ns) as the real measured value of the path --- and the
calibration run on the AMD machine, which built 250 of 250 tier binaries and calibrated the build
cadence to 0.521~s per binary and the space requirement to roughly 457~kB per binary; from this
follows a storage requirement of roughly 240~GB for the complete build of 524288 tier binaries over
the four system permutations of the golden profile (two optimization levels times two SIMD
options). The smoke series --- 43 permutations in two sub-blocks: 27 Cartesian
SIMD$\times$layout$\times$allocator combinations plus 16 over node width, path compression,
internal search, and telemetry --- is available in Appendix~\ref{app:measurements}; its provenance
is the data state under which it was collected: seven varied axes of that count, of which SIMD and
telemetry are no longer organ axes in the canonical order. It approximates axis effects partly via
wall-clock proxies, and individual effects are instrumentation artefacts\footnote{Instrumentation
artefact (probe effect): the change of the observed behaviour by the measuring instrument
itself~\cite{gait1986probe}.} rather than organ effects; the caveats per series are listed in
Table~\ref{tab:le:limitierung} (Section~\ref{sec:measurements:limitations}). The apparatus-side
acceptance per package runs through the two-gate protocol (Section~\ref{sec:impl-contracts}). The
\emph{consequence}: RQ0 and RQ3 remain open as long as series~A has not run; all eighteen axes are
planned for variation (not measured).

\section{Limitations}\label{sec:limitations}
The following limits describe what a user cannot achieve with the system --- apart from the
natural and self-chosen limits of the machine on which they work; caveats of individual measurement
series are listed in the measurement appendix (Table~\ref{tab:le:limitierung}), implementation
states in the implementation chapter (Chapter~\ref{ch:impl}), and the implementation status of the
full build and the smoke series is stated in Section~\ref{sec:answers}. Two limits follow from the
store principle, the rest are derived from the architecture: machine-specific empirical
optimization is the essential trait of cache-aware methods --- ATLAS~\cite{whaley2001atlas} and
FFTW~\cite{frigo1998fftw} survey the machine before they choose the code ---, and cache-oblivious
methods~\cite{frigo1999cache} are the opposite pole without machine binding; an apparatus that
makes machine binding its principle inherits its costs. The module bodies are not a limit: the code
generator of the measurement path emits per permutation the eighteen fully qualified axis types
together with the umbrella include, the legacy path with the probe stub is separated from the
measurement path by a strict barrier, and the zero statements of the fixed-assigned pool families
are explicitly coded non-observation with reason (Section~\ref{sec:impl-pipeline};
Table~\ref{tab:le:limitierung}).
\begin{itemize}
  \item \textbf{Version change of an axis algorithm:} If the user changes the version of an axis
  algorithm, every tier binary whose overlay cut contains the file becomes a new identity --- the
  overlay member of the fingerprint hashes all sources of the cut, and the version lock demands the
  deliberate regeneration commit for every change to a recorded file
  (Section~\ref{sec:impl-contracts}); the store's stock of these binaries together with measurement
  and evaluation has to be regenerated, because the fingerprint of the CacheEngineBuilder is
  provenance, not a reuse criterion. What the user cannot do: exchange a single algorithm and keep
  the measurements of the remaining binaries.
  \item \textbf{Store per machine type and hardware configuration:} Measured values are keyed to
  the hardware identity of the measuring machine --- stock~2 of the store is addressed via the
  fingerprint of the measured binary and the hardware identity, the platform cell bracketed beside
  it (Section~\ref{sec:lager}) ---; the entire store has to be measured separately per machine type
  and hardware configuration. What the user cannot do: transfer measured values from one machine
  to another. That is the price of machine-specific optimization (ATLAS~\cite{whaley2001atlas},
  FFTW~\cite{frigo1998fftw}) and the reason for the three machine signatures of the system
  registry.
  \item \textbf{Coverage of the hardware counters per microarchitecture:} In the measurement path
  the counter collection is active --- the emitted measurement jobs set the PMC switch, only the
  test build builds without hardware access. Collected are four generic counters (L1, last-level,
  dTLB, and branch misses via \texttt{perf\_event\_open}\footnote{\texttt{perf\_event\_open}(2): the
  Linux system call that opens a hardware or software counter as a file descriptor; generic events
  are mapped by the kernel per microarchitecture, model-specific ones require a raw encoding.})
  and two model-bound raw counters (L2 demand misses and fills from the cache of a foreign CCX as
  coherence observables), the latter only where the encoding is cross-checked by measurement ---
  AMD Zen~5; without a catalogue entry the fields carry an explicitly coded
  non-observation\footnote{Explicitly coded non-observation: the missing value is a state of its own
  (n/a with reason), never a zero; statistics speaks of planned missing data whose missingness
  mechanism is known and which therefore introduces no bias~\cite{graham2006planned,graham2009missing}.}
  (n/a), never a zero, and for Intel there is no raw assignment until an Intel host provides the
  cross-check. The last-level event fails to open on Zen~5 (no generic mapping),
  RAPL\footnote{RAPL (Running Average Power Limit): the energy counters of the processor packages,
  readable under Linux via the powercap sysfs zone; the zone is often readable only with root
  privileges.} is best-effort and empty without zone read permission, and IPC is declared as a
  category, not collected. In the smoke series the hardware columns are not collected (n/a); the
  tabulated zeros of the appendix tables are presentation, not measurement
  (Table~\ref{tab:le:limitierung}, row~1). What the user cannot do: obtain L2 and coherence
  counters on models other than Zen~5 (Figure~\ref{fig:measurement-triangle}).
  \item \textbf{Hardware fit without automatic check:} The hardware detection is available in
  three functions: the acquisition (the ISA complex axis with the SPD/XMP declaration of the memory
  clocking and the planner probe, which asks the host for its counters per run), the compile-time
  factory (the tier stamps from the actual properties), and the runtime verdict; the machine
  declarations of the system registry are declarations with a preserved source
  (\texttt{dmidecode}). Not present are the wiring into runner and probe dock, the additive result
  write-back, the provenance round trip of the declarations, and the fit stamp. What the user
  cannot do: have the fit between machine and binary checked automatically
  (Sections~\ref{sec:hw-probe} and~\ref{sec:impl-system-axes}).
  \item \textbf{Output without sheet fan-out:} The result workbook arises unconditionally in
  memory; xlsx is the default, the CSV form its selectable counterpart, always generated from it
  --- not a single child document but an ordered tree of CSV files with one folder per sheet ---,
  and both together are selectable
  (Section~\ref{sec:eval-identity}). The workbook possesses a constant sheet key --- all rows in
  one sheet ---; the fan-out into one sheet per sub-axis permutation plus an info sheet is planned
  as a filter chain of the evaluation but not implemented; the CSV projection is a flat file, the
  tree form not implemented. What the user cannot do: obtain results per sub-axis permutation as
  separate sheets or obtain the CSV storage as a tree.
  \item \textbf{Permutation space and budget:} A complete build of the golden reference catalogue
  costs, per system permutation, roughly 224~GB and roughly 24~hours according to the catalogue's
  estimate, roughly 457~kB and 0.521~s per tier binary according to the calibration
  (Section~\ref{sec:answers}); the planner computes the cardinality per profile and offers budget
  inputs in the subcommand \texttt{simulate} (bytes per binary, store budget, measurement subset,
  candidates). What the user cannot do: materialize the full product space of all axis catalogues
  --- they choose catalogues or subsets, and the number of repetitions per measurement remains a
  trade-off between effort and confidence~\cite{kalibera2013rigorous}.
  \item \textbf{No runtime switch between tier binaries:} The hybrid stage is implemented as
  classification, dock layer, XML parser, and binary proxy
  (Section~\ref{sec:impl-hybrid-lager}); the \texttt{.so} export of the hybrid stage, the eviction,
  and the router are not implemented --- reroute and router presuppose measurement curves. What the
  user cannot do: operate a hybrid tier that switches between preloaded binaries at run time; the
  state handover on switching --- the memento reconciliation with the two release artefacts single
  and hybrid in the store --- is likewise designed, not implemented.
  \item \textbf{Operating-system regime:} The system axis \texttt{operating\_system} is declared
  with \texttt{linux}, \texttt{windows}, and \texttt{macos} and possesses one probe each;
  measurement takes place under root Linux with full counter access, the second regime (immutable
  operating system) is designed, not implemented (Section~\ref{sec:eval-platforms}), and a Windows
  counter driver is not available (Table~\ref{tab:le:limitierung}, row~1). What the user cannot
  do: collect measurement series under Windows or macOS or without perf permissions
  (\texttt{perf\_event\_paranoid}, RAPL zones).
  \item \textbf{No isolated re-measurement of a single publication:} An operating mode that offers
  the design constituents of exactly one publication as a closed axis set is not implemented; per
  publication there exist the profile XML (33) and the replacement merge, not the isolation against
  all remaining axes. What the user cannot do: measure a publication in isolation, faithfully to
  the original (Section~\ref{sec:outlook}).
  \item \textbf{No automatic selection of the best configuration:} The heuristics-driven full
  automation (classifier and binary packager) is not implemented; present are the
  measurement-driven increment of the binary selection (ranking per metric from the evaluation
  output) and the runtime-dynamic cache-line policy. What the user cannot do: have an axis
  configuration recommended automatically from load-profile and hardware features; the
  recommendation remains a manual derivation from the evaluation (Section~\ref{sec:outlook}).
\end{itemize}

\section{Outlook}\label{sec:outlook}
In the short term the following steps are due; implemented parts appear with their open remainder,
designed parts with their implementation status:
\begin{enumerate}[label=(\arabic*)]
  \item the complete build according to the declared scope (Section~\ref{sec:eval-build-scope}) and
  the delivery measurement over the 320-identity catalogue;
  \item the end-to-end run of the three mandatory measurement series on the registered measurement
  machines together with the generation of the measurement diagrams;
  \item the full build-out of the store: the binaries \emph{with} and \emph{without} measurement
  sensors as two stored release artefacts and the synthesized function curves as files ---
  measurement is always switched on, in the store only what is missing is re-measured
  (Section~\ref{sec:lager}); the build-out is not implemented;
  \item the hardware detection: wiring into runner and probe dock, result write-back, provenance
  round trip, and fit stamp;
  \item the sheet fan-out of the xlsx workbook per sub-axis permutation plus info sheet --- the
  workbook writer is implemented, the fan-out is not implemented;
  \item the cross-check of the Intel raw counters on prod2, so that L2 and coherence counters are
  collected there as well;
  \item the second operating-system regime (immutable operating system) and the Windows counter
  driver;
  \item the paper-isolation mode from RQ2.c: an operating mode that offers the design constituents
  of exactly one publication as a closed axis set in isolation, in order to measure them --- with
  measurement and system axes remaining free-standing --- faithfully to the original;
  \item the profile connection of the organ meta-meta axis \texttt{disk\_io} and the permutation
  build over the scheduling sub-axis;
  \item the full permutation of the workloads against all living-being binaries (RQ2);
  \item the evaluation of the H2 criterion after the run of series~A\@.
\end{enumerate}

\paragraph{Heuristic automation and binary delivery.} The mechanism demanded by the thesis
assignment (sub-task~7), which automatically derives the best configuration from the measurement
results, is not implemented in its heuristics-driven \emph{full} automation and is a subject of
future work; a first increment of the automatic selection and binary shipping is available as a
standalone tool (\texttt{best\_binary\_selector}) that ranks and delivers the best permutation per
metric from the evaluation output --- the xlsx workbook is the output and the default, the CSV form
its ordered tree of CSV files generated from it (one folder per sheet, no single child document),
and the tool reads the CSV projection, which the seam writes as one flat file (the tree form is
planned but not implemented). Besides the selection (commitment~1), the same tool strand
comprises the runtime-dynamic cache-line adaptation (commitment~2), which is pertinent to the
cache-line awareness of the thesis: from an evaluated profile aggregate the policy selector derives
the cache-line policy. The full build-out presupposes three components; the second is available,
and toward the first as well as the third the tool provides a first, purely measurement-driven
increment: (i) a \emph{heuristic/ML classifier} that maps from the load-profile and axis features
(operation mix, key characteristics, hardware identifier) to the recommended axis configuration;
(ii) the implemented \emph{XML parser} (import and export including the profile loader) for the
heuristic policies stored as an XML load-profile result and reusable per architecture focus; and
(iii) a \emph{binary packager} that delivers the overall algorithm selected in this way as a
standalone binary behind the unified search-algorithm interface. A prerequisite for all three is
that the end-to-end measurement series are fully available --- only the robust data basis across
all permutations and load profiles makes a measurement-driven, automatically controlled selection
meaningful; without this automation the configuration recommendations of the evaluation scheme
from Section~\ref{sec:eval-results} are derived manually once the data are available.

\medskip
In the medium to long term: the measurement of the PIM allocator family on PIM
hardware\footnote{PIM (Processing-in-Memory): memory devices with computing units of their own,
such as UPMEM DPUs or Samsung's HBM-PIM\@; a PIM allocator requires such hardware.} --- the
allocator A16 is integrated as an axis building block and falls back to the portable aligned
allocation without PIM hardware~\cite{pimmalloc2026} ---; the formal verification of selected
permutation classes in the style of StarMalloc~\cite{reitz2024starmalloc}\footnote{StarMalloc is
verified in F\textsuperscript{$\ast$} with the framework Steel, a separation logic for concurrent
programs; the verification of an axis building block in this style is outlook.} --- the allocator
family A13 is integrated, the verification itself is open ---; the transfer of the anatomy
principle to accelerator memory hierarchies (GPU), without a commitment of cluster availability;
and the hybrid components live ranking, state handover, and measurement-error interpolation
(Section~\ref{sec:heuristics-modes}). The four-repository architecture (Section~\ref{sec:repos})
--- cache engine, probe, super-repository, manuscript --- allows further probes to be brought in
analogously to PRT-ART without changing the building-block library: the probe contract of the
build configuration checks the code path and the registry path against each other, so that both
halves of the integration point to the same place. Beyond the tree-like search structures treated
here, the anatomy principle is extensible via higher-level abstractions to the further structure
classes from Section~\ref{sec:overview} --- hash-based, spatial, and flat structures ---, so that
the same measurement and permutation apparatus can be applied to new classes without a rebuild;
herein lies an essential lever for future experimental progress.
